% ============================================================
%  teacher_guide.tex — The Secret Life of the Internet
%  Parent / Teacher Guide: discussion questions, extension
%  activities, and curriculum alignment suggestions.
% ============================================================

\namedchapter{Parent \& Teacher Guide}
\addcontentsline{toc}{chapter}{Parent \& Teacher Guide}

% REVIEW NOTE: Needs Minor Revision — improve responsive readability, accessibility tagging, and figure scaling before publication.
\begin{tcolorbox}[enhanced,colback=SunYellow!12,colframe=SunYellow!70!black,arc=6pt,boxrule=1pt,left=10pt,right=10pt,top=8pt,bottom=8pt]
\textbf{Review support note (Chapters 4--6):} These chapters are cognitively dense for many learners. Add recap boxes, step-by-step diagrams, or guided discussion prompts every 2--3 pages to maintain comprehension.\\
\textbf{Diagram readability note:} For packet-routing, server maps, and undersea cable figures, zoom in on digital displays or print at a larger scale for classroom use.
\end{tcolorbox}

\begin{tcolorbox}[enhanced,colback=LightGray,colframe=InternetBlue!50,arc=6pt,boxrule=1pt,
  left=10pt,right=10pt,top=10pt,bottom=10pt]
\textbf{How to use this guide:}
Each chapter entry below includes discussion questions to explore with your child or class,
extension activities for deeper learning, and notes on which curriculum areas the content supports.
The book is designed to be read one chapter at a time, with plenty of space for conversation between chapters.
\end{tcolorbox}

\vspace{12pt}

% ============================================================
\section*{Chapter 1: What Is the Internet?}
% ============================================================

\subsection*{Discussion Questions}
\begin{enumerate}
  \item How many devices in our home are connected to the internet? Can we count them together?
  \item What would change in your daily life if the internet stopped working for a week?
  \item The book says the internet is like a road system. Can you think of another analogy?
\end{enumerate}

\subsection*{Extension Activities}
\begin{itemize}
  \item \textbf{Home network audit:} Walk through the house and list every internet-connected device.
  \item \textbf{Historical timeline:} Research when the internet was invented (1969 ARPANET) and create a timeline of key milestones.
  \item \textbf{Human network game:} Stand in a circle, pass a ball (the ``packet'') from person to person — each person is a node in the network.
\end{itemize}

\subsection*{Curriculum Links}
Science \& Technology, Computing, Geography (global connections).

\sectionrule

% ============================================================
\section*{Chapter 2: Websites and Browsers}
% ============================================================

\subsection*{Discussion Questions}
\begin{enumerate}
  \item What browser do we use at home? Have you ever tried a different one?
  \item Why do websites have different endings like .com, .org, .ca, .uk?
  \item What do you think happens if two people try to register the same website name?
\end{enumerate}

\subsection*{Extension Activities}
\begin{itemize}
  \item \textbf{View page source:} Right-click any web page and select ``View Source'' to see real HTML code.
  \item \textbf{Build a mini page:} Use a simple text editor to create a basic HTML page with a heading, paragraph, and image.
  \item \textbf{URL scavenger hunt:} Give children 5 URLs to visit and describe what kind of ``building'' each website would be.
\end{itemize}

\subsection*{Curriculum Links}
Computing (programming concepts), English (reading for information).

\sectionrule

% ============================================================
\section*{Chapter 3: Wi-Fi and Routers}
% ============================================================

\subsection*{Discussion Questions}
\begin{enumerate}
  \item Why does Wi-Fi sometimes work better in one room than another?
  \item What is the difference between Wi-Fi and mobile data (4G/5G)?
  \item Can you think of other things that use radio waves besides Wi-Fi?
\end{enumerate}

\subsection*{Extension Activities}
\begin{itemize}
  \item \textbf{Signal strength walk:} Use a phone to check Wi-Fi signal strength in different rooms and graph the results.
  \item \textbf{Wave experiment:} Drop objects in water to observe how waves travel and weaken with distance.
  \item \textbf{Router investigation:} Look at your home router — how many lights does it have? What do they mean?
\end{itemize}

\subsection*{Curriculum Links}
Science (waves, energy transfer), Mathematics (measurement, graphing).

\sectionrule

% ============================================================
\section*{Chapter 4: Servers and Data Centres}
% ============================================================

\subsection*{Discussion Questions}
\begin{enumerate}
  \item Why do data centres need so much cooling?
  \item What would happen if a data centre lost power and had no backup?
  \item Why do big companies build data centres in many different countries?
\end{enumerate}

\subsection*{Extension Activities}
\begin{itemize}
  \item \textbf{Virtual tour:} Watch a video tour of a real data centre (Google and Microsoft publish these).
  \item \textbf{Request timer:} Use a stopwatch to time how long it takes different websites to load — discuss why some are faster.
  \item \textbf{Energy awareness:} Research how much electricity a data centre uses and compare to household usage.
\end{itemize}

\subsection*{Curriculum Links}
Science (energy, heat), Geography (global infrastructure), PSHE (environmental responsibility).

\sectionrule

% ============================================================
\section*{Chapter 5: Undersea Cables}
% ============================================================

\subsection*{Discussion Questions}
\begin{enumerate}
  \item Why are cables better than satellites for carrying most internet traffic?
  \item How do engineers fix a cable that breaks at the bottom of the ocean?
  \item The first transatlantic cable took 16 hours to send 98 words. How long does a message take today?
\end{enumerate}

\subsection*{Extension Activities}
\begin{itemize}
  \item \textbf{Submarine cable map:} Visit \url{https://submarinecablemap.com} and trace a route from your country to another continent.
  \item \textbf{Light experiment:} Shine a torch through a clear straw or water stream to demonstrate total internal reflection (how fibre optics work).
  \item \textbf{Scale model:} Use string and a world map to create a model of undersea cable routes.
\end{itemize}

\subsection*{Curriculum Links}
Science (light, total internal reflection), History (communication technology), Geography (oceans, continents).

\sectionrule

% ============================================================
\section*{Chapter 6: How Messages and Videos Travel}
% ============================================================

\subsection*{Discussion Questions}
\begin{enumerate}
  \item Why is it useful to split a message into packets instead of sending it whole?
  \item What happens when one packet gets lost? How does the system recover?
  \item Why does a video sometimes buffer when your internet is slow?
\end{enumerate}

\subsection*{Extension Activities}
\begin{itemize}
  \item \textbf{Packet relay race:} Write a sentence on paper, cut it into strips (packets), number them, hand them to different runners — reassemble at the finish line.
  \item \textbf{Compression demo:} Compare file sizes of the same photo at different quality levels.
  \item \textbf{Traceroute:} On a computer, run \texttt{traceroute} (Mac/Linux) or \texttt{tracert} (Windows) to see the hops a packet takes to reach a website.
\end{itemize}

\subsection*{Curriculum Links}
Mathematics (sequencing, ordering), Computing (protocols, algorithms), PE (relay games).

\sectionrule

% ============================================================
\section*{Chapter 7: Search Engines}
% ============================================================

\subsection*{Discussion Questions}
\begin{enumerate}
  \item How does a search engine decide which results to show first?
  \item Why might two people searching the same words get slightly different results?
  \item How can you tell if a website is trustworthy?
\end{enumerate}

\subsection*{Extension Activities}
\begin{itemize}
  \item \textbf{Search challenge:} Give children an obscure fact to find — who can find a reliable source fastest?
  \item \textbf{Compare engines:} Try the same search on Google, Bing, and DuckDuckGo — how do results differ?
  \item \textbf{Fact-check exercise:} Present a surprising claim and ask children to verify it using at least two sources.
\end{itemize}

\subsection*{Curriculum Links}
English (research skills, critical reading), PSHE (media literacy), Computing (information retrieval).

\sectionrule

% ============================================================
\section*{Chapter 8: The Cloud}
% ============================================================

\subsection*{Discussion Questions}
\begin{enumerate}
  \item What files or photos do we have stored in the cloud right now?
  \item What are the advantages and disadvantages of keeping files in the cloud vs.\ on your device?
  \item How does the cloud help people work together from different countries?
\end{enumerate}

\subsection*{Extension Activities}
\begin{itemize}
  \item \textbf{Collaborative document:} Create a shared document and have family members edit it from different devices simultaneously.
  \item \textbf{Offline test:} Turn off Wi-Fi and see which apps and files still work — discuss why.
  \item \textbf{Storage calculator:} Estimate how many photos, songs, or videos fit in 1 GB of cloud storage.
\end{itemize}

\subsection*{Curriculum Links}
Computing (storage, collaboration tools), Mathematics (data sizes, estimation), PSHE (digital responsibility).

\sectionrule

% ============================================================
\section*{Chapter 9: Stay Safe and Curious}
% ============================================================

\subsection*{Discussion Questions}
\begin{enumerate}
  \item What personal information should you never share online without asking a trusted adult?
  \item How would you respond if someone online asked you to keep a secret from your parents?
  \item Why is being kind online just as important as being kind in person?
\end{enumerate}

\subsection*{Extension Activities}
\begin{itemize}
  \item \textbf{Password strength test:} Create passwords of different lengths and discuss which would be hardest to guess.
  \item \textbf{Phishing spotting:} Show examples of suspicious emails (with identifying info removed) and practise spotting red flags.
  \item \textbf{Screen time diary:} Track screen time for a week, then discuss the balance between online and offline activities.
  \item \textbf{Digital citizenship poster:} Create a poster of the ``Digital Explorer's Code of Honour'' for the classroom or fridge.
\end{itemize}

\subsection*{Curriculum Links}
PSHE (online safety, well-being, relationships), Computing (digital citizenship), Art (poster design).

\sectionrule

% ============================================================
\section*{Suggested Reading Schedule}
% ============================================================

\begin{tcolorbox}[enhanced,colback=SunYellow!10,colframe=SunYellow!70!black,arc=6pt,boxrule=1pt,
  left=10pt,right=10pt,top=8pt,bottom=8pt]

\begin{center}
\small
\renewcommand{\arraystretch}{1.2}
\begin{tabularx}{\linewidth}{C{0.12\linewidth} C{0.24\linewidth} Y}
\textbf{Week} & \textbf{Chapter} & \textbf{Focus} \\
\hline
1 & Chapter 1 & What is the internet? Drawing networks. \\
2 & Chapters 2--3 & Websites, browsers, Wi-Fi, routers. \\
3 & Chapters 4--5 & Servers, data centres, undersea cables. \\
4 & Chapters 6--7 & Packets, search engines. \\
5 & Chapters 8--9 & The cloud, online safety. Quiz \& review. \\
\end{tabularx}
\end{center}

This five-week schedule works well for classroom use or family read-aloud sessions,
allowing time for discussion and extension activities between chapters.
\end{tcolorbox}
